
    \section{Matching}
        \subsection{Project Selection}
        Suppose we have the following: we have $N$ projects and $M$ machines. The $i$-th machine costs $c_i$. The $i$-th project gives $v_i$ money. Some projects requires some machines to be done, and if more than one machine needs the same project, they can be shared. You want to maximize the revenue generated by the projects.

        First, let's make $X = \sum_{x \in S} v_i$ the sum of all projects. By doing this, we are assuming we have the money for doing all the projects, and now what we will do is return the money of the projects we didn't manage to finish.

        The way this modeling works is: This graph is bipartite when looking at the costs: some are positive, and some are negative. Let's link the source with every project, using its cost as the flow passing by its edge, and let's link the machines on the sink by its price to be used. After it, let's link the projects to its associated machines, with $\infty$. This is the way we will be constructing the graph. It has $N+M+2$ nodes, $N$ being the projects, $M$ being the machines and $2$ more for the source and sink nodes.

        Now, to reconstruct the solution, we will do a Minimum Cut on the graph. If an edge between the source and a project $i$ is cut, then we could not finish it, and we need to return its money. If an edge between a machine $i$ and the sink is cut, then we need to buy this machine $i$. Since it is impossible to cut the $\infty$ edges between projects and machines, the edges presented here will be the only ones to be cut by the algorithm.
        
        O problema é o seguinte, dado:
            \begin{itemize}
                \item um conjunto de ferramentas, cada uma custando $f_i$.
                \item um conjunto de projetos, cada um entregando um lucro $p_i$.
                \item um conjunto de requisitos $r_{i,j}$ que são necessários para executar o projeto $i$.
            \end{itemize}

        Queremos maximizar nosso lucro. Sendo lucro $ = \sum p_j - \sum f_k$, sendo $j$ os projetos escolhidos e $k$ os projetos requisitados para pegar os $j$ projetos.

        A solução é a seguinte: Modelaremos o nosso problema como um corte mínimo em um grafo. Ligando todos os projetos ao $sink$ com a capacidade sendo $p_i$ e todas as ferramentas ao $source$ com capacidade sendo $f_i$. Entre eles colocaremos uma capacidade infinita, conectando os projetos com suas respectivas ferramentas.

        Assim, nossa resposta será $\sum p_i - \text{MinCut}$.
        
        \subsection {HopCroft Karp}
            Compĺexity for sparse graphs $O(E log(V))$
            Compĺexity for dense graphs $O(E sqrt(V))$
            \inputwithhash{./solutions/Graph/non_weighted_bipartite_matching.cpp}
        \subsection{Bipartide Weighted Hungarian Method}
        Compĺexity $O(V^3)$
        \inputwithhash{./solutions/Graph/hungarian.cpp}
        \subsection{General Graph (Edmonds Blossom)}
        Complexity $O(VE^2)$
        \inputwithhash{./solutions/Graph/MatchGeneral.cpp}
